\documentclass[11pt]{article}

\usepackage[margin=0.9in]{geometry}
\usepackage{amsmath,amssymb}
\usepackage{booktabs}
\usepackage{enumitem}
\usepackage{listings}
\usepackage{xcolor}
\usepackage{hyperref}

\definecolor{backcolour}{rgb}{0.95,0.95,0.92}
\definecolor{codeblue}{rgb}{0.05,0.15,0.55}
\definecolor{codegreen}{rgb}{0.0,0.45,0.0}

\lstset{
    backgroundcolor=\color{backcolour},
    basicstyle=\ttfamily\small,
    keywordstyle=\color{codeblue}\bfseries,
    commentstyle=\color{codegreen},
    stringstyle=\color{purple},
    breaklines=true,
    frame=single,
    showstringspaces=false,
    tabsize=4
}

\title{CPSC 250 \\ Basic Plotting of Data}
\author{Edward Brash}
\date{}

\begin{document}

\maketitle

\section{Motivation}

We have already seen examples where plotting data is useful.

In this lecture, we develop a base plotting example using \texttt{matplotlib}.

The goal is to create a plot with:

\begin{enumerate}
    \item a title,
    \item axis labels,
    \item data points with error bars,
    \item a fit line,
    \item a legend,
    \item appropriate axis limits,
    \item and possibly logarithmic axes or an error band.
\end{enumerate}

\section{Input Data Format}

Suppose we have data stored in a CSV file.

The file has columns:

\[
x,\quad y,\quad \Delta x,\quad \Delta y
\]

where:

\begin{itemize}
    \item \(x\) is the independent variable,
    \item \(y\) is the dependent variable,
    \item \(\Delta x\) is the uncertainty in \(x\),
    \item \(\Delta y\) is the uncertainty in \(y\).
\end{itemize}

A schematic CSV file might look like:

\begin{lstlisting}
XLabel,YLabel,dXLabel,dYLabel
1.0,3.0,0.01,0.03
2.0,3.2,0.02,0.02
...
\end{lstlisting}

\section{What We Want to Plot}

A good data plot should show:

\begin{itemize}
    \item the data points,
    \item both x and y error bars,
    \item the best-fit curve,
    \item a legend describing the data and fit,
    \item the x-axis label,
    \item the y-axis label,
    \item a title.
\end{itemize}

\section{Reading the Data}

We can use Python's \texttt{csv} module to read the data.

\begin{lstlisting}[language=Python]
import csv

def read_data(filename):

    x = []
    y = []
    dx = []
    dy = []

    with open(filename, "r") as file:

        reader = csv.reader(file)

        headers = next(reader, None)

        for row in reader:

            x.append(float(row[0]))
            y.append(float(row[1]))
            dx.append(float(row[2]))
            dy.append(float(row[3]))

    return headers, x, y, dx, dy
\end{lstlisting}

\section{The Header Row}

The line:

\begin{lstlisting}[language=Python]
headers = next(reader, None)
\end{lstlisting}

reads the first row of the file.

This row contains labels rather than data.

For example:

\begin{lstlisting}
XLabel,YLabel,dXLabel,dYLabel
\end{lstlisting}

We can use these labels later for axis labels.

\section{Creating the Basic Plot}

We import:

\begin{lstlisting}[language=Python]
import matplotlib.pyplot as plt
\end{lstlisting}

Then:

\begin{lstlisting}[language=Python]
filename = "testdata.csv"

headers, xi, yi, dxi, dyi = read_data(filename)

plt.errorbar(xi, yi, xerr=dxi, yerr=dyi, fmt="o")

plt.title("Basic Plotting Example")
plt.xlabel(headers[0])
plt.ylabel(headers[1])

plt.show()
\end{lstlisting}

\section{The \texttt{errorbar()} Function}

The function:

\begin{lstlisting}[language=Python]
plt.errorbar()
\end{lstlisting}

plots points with error bars.

In this example:

\begin{lstlisting}[language=Python]
plt.errorbar(xi, yi, xerr=dxi, yerr=dyi, fmt="o")
\end{lstlisting}

means:

\begin{itemize}
    \item use \texttt{xi} as x-values,
    \item use \texttt{yi} as y-values,
    \item use \texttt{dxi} as x uncertainties,
    \item use \texttt{dyi} as y uncertainties,
    \item use circular markers.
\end{itemize}

\section{Axis Limits}

Matplotlib will choose default axis limits automatically.

However, sometimes we should set the limits ourselves.

For example:

\begin{lstlisting}[language=Python]
plt.ylim(0, 0)
\end{lstlisting}

is a special case: Python interprets this as a request to choose the upper limit automatically while forcing the lower limit to zero.

This can be useful when the physical quantity cannot be negative.

\section{Which Plot Requirements Are Already Satisfied?}

At this stage, we have:

\begin{itemize}
    \item a title,
    \item x-axis label,
    \item y-axis label,
    \item data points with x and y error bars,
    \item appropriate automatic axis behavior.
\end{itemize}

We have not yet added:

\begin{itemize}
    \item a fit curve,
    \item a legend,
    \item logarithmic axes,
    \item an error band.
\end{itemize}

\section{Fitting}

Fitting is a very large topic.

In this course, we use a simple practical approach.

The steps are:

\begin{enumerate}
    \item Choose a fitting function.
    \item Call a fitting routine.
    \item Extract the best-fit parameters.
    \item Plot the fit.
\end{enumerate}

\section{Choosing a Fit Function}

Suppose the data appear roughly quadratic.

Then we might choose:

\[
y = ax^2 + bx + c
\]

In Python:

\begin{lstlisting}[language=Python]
def fitfunction(x, *param):

    return param[0]*x*x + param[1]*x + param[2]
\end{lstlisting}

The star in:

\begin{lstlisting}[language=Python]
*param
\end{lstlisting}

allows the function to accept a variable number of parameters.

Here:

\begin{itemize}
    \item \texttt{param[0]} is \(a\),
    \item \texttt{param[1]} is \(b\),
    \item \texttt{param[2]} is \(c\).
\end{itemize}

\section{Calling \texttt{curve\_fit}}

We import:

\begin{lstlisting}[language=Python]
from scipy.optimize import curve_fit
\end{lstlisting}

Then we call:

\begin{lstlisting}[language=Python]
init_vals = [0.0 for x in range(3)]

popt, pcov = curve_fit(
    fitfunction,
    xi,
    yi,
    p0=init_vals,
    sigma=dyi,
    absolute_sigma=True
)
\end{lstlisting}

Here:

\begin{itemize}
    \item \texttt{popt} contains the optimal fit parameters,
    \item \texttt{pcov} contains the covariance matrix.
\end{itemize}

\section{Parameter Uncertainties}

The uncertainties in the fitted parameters come from the diagonal of the covariance matrix:

\begin{lstlisting}[language=Python]
import numpy as np

perr = np.sqrt(np.diag(pcov))
\end{lstlisting}

Then:

\begin{lstlisting}[language=Python]
a = popt[0]
b = popt[1]
c = popt[2]

da = perr[0]
db = perr[1]
dc = perr[2]
\end{lstlisting}

This gives:

\[
a \pm \Delta a
\]

\[
b \pm \Delta b
\]

\[
c \pm \Delta c
\]

\section{Plotting the Fit}

To plot a smooth curve, create many x-values:

\begin{lstlisting}[language=Python]
xlow = min(xi) - 2
xhigh = max(xi) + 2

xfit = np.linspace(xlow, xhigh, 100)
yfit = fitfunction(xfit, *popt)

plt.plot(xfit, yfit, "--")
\end{lstlisting}

The line:

\begin{lstlisting}[language=Python]
fitfunction(xfit, *popt)
\end{lstlisting}

evaluates the fit function using the fitted parameters.

\section{Key Takeaways}

\begin{itemize}
    \item CSV files are useful for storing data.
    \item A good plot includes labels, a title, and a legend.
    \item Error bars communicate measurement uncertainty.
    \item \texttt{plt.errorbar()} plots data with uncertainties.
    \item \texttt{curve\_fit()} performs numerical fitting.
    \item The covariance matrix provides parameter uncertainties.
    \item Smooth fit curves are created using many x-values.
\end{itemize}

\end{document}
